% The "MLJ Contribution Information Sheet" that the Springer journal
% "Machine Learning" requires with every submission: a 1-2 page document
% answering four fixed questions (main claim and its importance, evidence,
% closest related work, prior publication). Submissions with an incomplete
% sheet are returned without review. It is uploaded as a separate file
% next to the manuscript. The numbers come from the same generated
% fragments as the article, so the two cannot disagree.
\documentclass[10pt,a4paper]{article}
\usepackage[margin=2cm]{geometry}
\usepackage[T1]{fontenc}
\usepackage[utf8]{inputenc}
\usepackage{microtype}
\usepackage[natbibapa]{apacite}
\usepackage[hidelinks]{hyperref}

\newcommand{\useresults}[1]{\IfFileExists{#1.gen.tex}{\input{#1.gen.tex}}{}}
\useresults{adaptive-2wiki}
\useresults{adaptive-hotpotqa}
\useresults{adaptive-musique}
\useresults{adaptive-nq}
\useresults{adaptive-triviaqa}
\useresults{bm25-2wiki}
\useresults{bm25-hotpotqa}
\useresults{bm25-musique}
\useresults{bm25-nq}
\useresults{bm25-triviaqa}
\useresults{cascade-2wiki}
\useresults{cascade-hotpotqa}
\useresults{cascade-hybrid-2wiki}
\useresults{cascade-hybrid-hotpotqa}
\useresults{cascade-hybrid-musique}
\useresults{cascade-hybrid-nq}
\useresults{cascade-hybrid-triviaqa}
\useresults{cascade-llama-2wiki}
\useresults{cascade-llama-hotpotqa}
\useresults{cascade-llama-musique}
\useresults{cascade-llama-nq}
\useresults{cascade-llama-triviaqa}
\useresults{cascade-lp-2wiki}
\useresults{cascade-lp-hotpotqa}
\useresults{cascade-lp-musique}
\useresults{cascade-lp-nq}
\useresults{cascade-lp-triviaqa}
\useresults{cascade-musique}
\useresults{cascade-naive-2wiki}
\useresults{cascade-naive-hotpotqa}
\useresults{cascade-naive-musique}
\useresults{cascade-naive-nq}
\useresults{cascade-naive-triviaqa}
\useresults{cascade-nq}
\useresults{cascade-rr-2wiki}
\useresults{cascade-rr-hotpotqa}
\useresults{cascade-rr-musique}
\useresults{cascade-rr-nq}
\useresults{cascade-rr-triviaqa}
\useresults{cascade-triviaqa}
\useresults{closedbook-2wiki}
\useresults{closedbook-hotpotqa}
\useresults{closedbook-llama-2wiki}
\useresults{closedbook-llama-hotpotqa}
\useresults{closedbook-llama-musique}
\useresults{closedbook-llama-nq}
\useresults{closedbook-llama-triviaqa}
\useresults{closedbook-musique}
\useresults{closedbook-nq}
\useresults{closedbook-triviaqa}
\useresults{cmp-cascade-cascade-lp-2wiki}
\useresults{cmp-cascade-cascade-lp-hotpotqa}
\useresults{cmp-cascade-cascade-lp-musique}
\useresults{cmp-cascade-cascade-lp-nq}
\useresults{cmp-cascade-cascade-lp-triviaqa}
\useresults{cmp-cascade-hybrid-cascade-2wiki}
\useresults{cmp-cascade-hybrid-cascade-hotpotqa}
\useresults{cmp-cascade-hybrid-cascade-musique}
\useresults{cmp-cascade-hybrid-cascade-nq}
\useresults{cmp-cascade-hybrid-cascade-triviaqa}
\useresults{cmp-cascade-naive-cascade-2wiki}
\useresults{cmp-cascade-naive-cascade-hotpotqa}
\useresults{cmp-cascade-naive-cascade-musique}
\useresults{cmp-cascade-naive-cascade-nq}
\useresults{cmp-cascade-naive-cascade-triviaqa}
\useresults{cmp-cascade-rr-cascade-2wiki}
\useresults{cmp-cascade-rr-cascade-hotpotqa}
\useresults{cmp-cascade-rr-cascade-musique}
\useresults{cmp-cascade-rr-cascade-nq}
\useresults{cmp-cascade-rr-cascade-triviaqa}
\useresults{cmp-closedbook-naive-2wiki}
\useresults{cmp-closedbook-naive-hotpotqa}
\useresults{cmp-closedbook-naive-llama-2wiki}
\useresults{cmp-closedbook-naive-llama-hotpotqa}
\useresults{cmp-closedbook-naive-llama-musique}
\useresults{cmp-closedbook-naive-llama-nq}
\useresults{cmp-closedbook-naive-llama-triviaqa}
\useresults{cmp-closedbook-naive-musique}
\useresults{cmp-closedbook-naive-nq}
\useresults{cmp-closedbook-naive-triviaqa}
\useresults{cmp-fused-cascade-2wiki}
\useresults{cmp-fused-cascade-hotpotqa}
\useresults{cmp-fused-cascade-llama-2wiki}
\useresults{cmp-fused-cascade-llama-hotpotqa}
\useresults{cmp-fused-cascade-llama-musique}
\useresults{cmp-fused-cascade-llama-nq}
\useresults{cmp-fused-cascade-llama-triviaqa}
\useresults{cmp-fused-cascade-musique}
\useresults{cmp-fused-cascade-nq}
\useresults{cmp-fused-cascade-triviaqa}
\useresults{cmp-ircot-adaptive-2wiki}
\useresults{cmp-ircot-adaptive-hotpotqa}
\useresults{cmp-ircot-adaptive-musique}
\useresults{cmp-ircot-adaptive-nq}
\useresults{cmp-ircot-adaptive-triviaqa}
\useresults{cmp-ircot-ircot1-2wiki}
\useresults{cmp-ircot-ircot1-hotpotqa}
\useresults{cmp-ircot-ircot1-musique}
\useresults{cmp-ircot-ircot1-nq}
\useresults{cmp-ircot-ircot1-triviaqa}
\useresults{cmp-ircot1-cascade-2wiki}
\useresults{cmp-ircot1-cascade-hotpotqa}
\useresults{cmp-ircot1-cascade-musique}
\useresults{cmp-ircot1-cascade-nq}
\useresults{cmp-ircot1-cascade-triviaqa}
\useresults{cmp-naive-adaptive-2wiki}
\useresults{cmp-naive-adaptive-hotpotqa}
\useresults{cmp-naive-adaptive-musique}
\useresults{cmp-naive-adaptive-nq}
\useresults{cmp-naive-adaptive-triviaqa}
\useresults{cmp-naive-bm25-2wiki}
\useresults{cmp-naive-bm25-hotpotqa}
\useresults{cmp-naive-bm25-musique}
\useresults{cmp-naive-bm25-nq}
\useresults{cmp-naive-bm25-triviaqa}
\useresults{cmp-naive-cascade-2wiki}
\useresults{cmp-naive-cascade-hotpotqa}
\useresults{cmp-naive-cascade-llama-2wiki}
\useresults{cmp-naive-cascade-llama-hotpotqa}
\useresults{cmp-naive-cascade-llama-musique}
\useresults{cmp-naive-cascade-llama-nq}
\useresults{cmp-naive-cascade-llama-triviaqa}
\useresults{cmp-naive-cascade-lp-2wiki}
\useresults{cmp-naive-cascade-lp-hotpotqa}
\useresults{cmp-naive-cascade-lp-musique}
\useresults{cmp-naive-cascade-lp-nq}
\useresults{cmp-naive-cascade-lp-triviaqa}
\useresults{cmp-naive-cascade-musique}
\useresults{cmp-naive-cascade-nq}
\useresults{cmp-naive-cascade-triviaqa}
\useresults{cmp-naive-hybrid-2wiki}
\useresults{cmp-naive-hybrid-hotpotqa}
\useresults{cmp-naive-hybrid-musique}
\useresults{cmp-naive-hybrid-nq}
\useresults{cmp-naive-hybrid-triviaqa}
\useresults{cmp-naive-hyde-2wiki}
\useresults{cmp-naive-hyde-hotpotqa}
\useresults{cmp-naive-hyde-musique}
\useresults{cmp-naive-hyde-nq}
\useresults{cmp-naive-hyde-triviaqa}
\useresults{cmp-naive-rerank-2wiki}
\useresults{cmp-naive-rerank-hotpotqa}
\useresults{cmp-naive-rerank-llama-2wiki}
\useresults{cmp-naive-rerank-llama-hotpotqa}
\useresults{cmp-naive-rerank-llama-musique}
\useresults{cmp-naive-rerank-llama-nq}
\useresults{cmp-naive-rerank-llama-triviaqa}
\useresults{cmp-naive-rerank-musique}
\useresults{cmp-naive-rerank-nq}
\useresults{cmp-naive-rerank-triviaqa}
\useresults{cmp-naive10-crag10-2wiki}
\useresults{cmp-naive10-crag10-hotpotqa}
\useresults{cmp-naive10-crag10-musique}
\useresults{cmp-naive10-ircot-2wiki}
\useresults{cmp-naive10-ircot-hotpotqa}
\useresults{cmp-naive10-ircot-musique}
\useresults{cmp-naive10-ircot1-2wiki}
\useresults{cmp-naive10-ircot1-hotpotqa}
\useresults{cmp-naive10-ircot1-musique}
\useresults{cmp-naive12-neighbour-2wiki}
\useresults{cmp-naive12-neighbour-hotpotqa}
\useresults{cmp-naive12-neighbour-musique}
\useresults{cmp-rerank-cascade-2wiki}
\useresults{cmp-rerank-cascade-hotpotqa}
\useresults{cmp-rerank-cascade-llama-2wiki}
\useresults{cmp-rerank-cascade-llama-hotpotqa}
\useresults{cmp-rerank-cascade-llama-musique}
\useresults{cmp-rerank-cascade-llama-nq}
\useresults{cmp-rerank-cascade-llama-triviaqa}
\useresults{cmp-rerank-cascade-lp-2wiki}
\useresults{cmp-rerank-cascade-lp-hotpotqa}
\useresults{cmp-rerank-cascade-lp-musique}
\useresults{cmp-rerank-cascade-lp-nq}
\useresults{cmp-rerank-cascade-lp-triviaqa}
\useresults{cmp-rerank-cascade-musique}
\useresults{cmp-rerank-cascade-nq}
\useresults{cmp-rerank-cascade-triviaqa}
\useresults{cmp-rerank-fused-2wiki}
\useresults{cmp-rerank-fused-hotpotqa}
\useresults{cmp-rerank-fused-llama-2wiki}
\useresults{cmp-rerank-fused-llama-hotpotqa}
\useresults{cmp-rerank-fused-llama-musique}
\useresults{cmp-rerank-fused-llama-nq}
\useresults{cmp-rerank-fused-llama-triviaqa}
\useresults{cmp-rerank-fused-musique}
\useresults{cmp-rerank-fused-nq}
\useresults{cmp-rerank-fused-triviaqa}
\useresults{cost}
\useresults{coverage-2wiki}
\useresults{coverage-hotpotqa}
\useresults{coverage-musique}
\useresults{crag-2wiki}
\useresults{crag-hotpotqa}
\useresults{crag-musique}
\useresults{crag-nq}
\useresults{crag-triviaqa}
\useresults{crag10-2wiki}
\useresults{crag10-hotpotqa}
\useresults{crag10-musique}
\useresults{diagnosis-2wiki}
\useresults{diagnosis-hotpotqa}
\useresults{diagnosis-musique}
\useresults{fused-2wiki}
\useresults{fused-hotpotqa}
\useresults{fused-llama-2wiki}
\useresults{fused-llama-hotpotqa}
\useresults{fused-llama-musique}
\useresults{fused-llama-nq}
\useresults{fused-llama-triviaqa}
\useresults{fused-musique}
\useresults{fused-nq}
\useresults{fused-triviaqa}
\useresults{hybrid-2wiki}
\useresults{hybrid-hotpotqa}
\useresults{hybrid-musique}
\useresults{hybrid-nq}
\useresults{hybrid-triviaqa}
\useresults{hyde-2wiki}
\useresults{hyde-hotpotqa}
\useresults{hyde-musique}
\useresults{hyde-nq}
\useresults{hyde-triviaqa}
\useresults{ircot-2wiki}
\useresults{ircot-hotpotqa}
\useresults{ircot-musique}
\useresults{ircot-nq}
\useresults{ircot-triviaqa}
\useresults{ircot1-2wiki}
\useresults{ircot1-hotpotqa}
\useresults{ircot1-musique}
\useresults{ircot1-nq}
\useresults{ircot1-triviaqa}
\useresults{naive-2wiki}
\useresults{naive-hotpotqa}
\useresults{naive-llama-2wiki}
\useresults{naive-llama-hotpotqa}
\useresults{naive-llama-musique}
\useresults{naive-llama-nq}
\useresults{naive-llama-triviaqa}
\useresults{naive-musique}
\useresults{naive-nq}
\useresults{naive-triviaqa}
\useresults{naive10-2wiki}
\useresults{naive10-hotpotqa}
\useresults{naive10-musique}
\useresults{naive12-2wiki}
\useresults{naive12-hotpotqa}
\useresults{naive12-musique}
\useresults{neighbour-2wiki}
\useresults{neighbour-hotpotqa}
\useresults{neighbour-musique}
\useresults{rerank-2wiki}
\useresults{rerank-hotpotqa}
\useresults{rerank-llama-2wiki}
\useresults{rerank-llama-hotpotqa}
\useresults{rerank-llama-musique}
\useresults{rerank-llama-nq}
\useresults{rerank-llama-triviaqa}
\useresults{rerank-musique}
\useresults{rerank-nq}
\useresults{rerank-triviaqa}
\useresults{result-llama}
\useresults{result-lp}
\useresults{result}
\useresults{sweep}


\title{MLJ Contribution Information Sheet}
\author{}
\date{}

\begin{document}
\maketitle
\vspace{-3em}
\noindent\textbf{Manuscript:} \emph{CascadeRAG: Abstention-Triggered
Escalation that Outperforms Costlier Retrieval-Augmented Generation
Architectures at a Fraction of their Cost}\\
\textbf{Authors:} Miko{\l}aj Semeniuk (corresponding author) and Agnieszka
Jastrz{\k{e}}bska, Warsaw University of Technology, Warsaw, Poland

\section*{1. What is the main claim of the paper? Why is this an important
contribution to the machine learning literature?}

\paragraph{Main claim.} Among the architectural refinements proposed for
retrieval-augmented generation (RAG), none is universally safe, and the
design that wins a controlled comparison is a cascade that spends extra
computation only where a free, posterior signal says it is needed. The
paper makes three concrete claims. (i)~Under one controlled stack, the
strongest single refinement, cross-encoder reranking, gains $+0.051$ Exact
Match on one question set and loses $0.100$ on another, a reversal that
single-benchmark evaluations do not reveal. (ii)~A per-question diagnosis
of \emph{why} retrieval fails (five causes) predicts which refinement pays
off on which question set. (iii)~CascadeRAG, a novel architecture built
from that diagnosis, fuses dense, lexical, and cross-encoder rankings and
escalates a question to a more expensive strategy only when the reader
itself, having seen the evidence, declines to answer. This abstention
trigger is a reject option with provably zero selective risk under Exact
Match, so escalation can never discard a correct answer. CascadeRAG wins
\ResultWins{} of \ResultPairs{} paired (baseline, question set)
comparisons, loses \ResultLosses{}, and ties \ResultTies{}, at 1.16--1.58
generator calls per question against 2.3--3.0 for iterative and corrective
methods.

\paragraph{Why it matters.} RAG is now the standard way of grounding large
language models in a document collection, yet its architectures are
compared in the literature as point scores on incompatible stacks. The
paper supplies two things the machine learning literature lacks here: a
research-methodology contribution, namely a comparison protocol with
paired significance tests under multiplicity control, budget-matched
controls, abstention accounted for separately from error, and a measured
corpus-imposed recall ceiling, under which most published-style gains
shrink and one reverses; and a learning-systems contribution that connects
RAG design to selective prediction, showing that a training-free,
posterior deferral signal dominates both prior routing and
confidence-threshold deferral. Two documented negative results
(neighbour-chunk expansion adds nothing once context size is controlled;
answer-token confidence adds nothing over abstention) save practitioners
from two widely assumed improvements.

\section*{2. What is the evidence you provide to support your claim? Be
precise.}

Thirteen RAG architectures are evaluated on a single stack: one
21-million-passage Wikipedia corpus, one embedding model, one vector
database, one generator (Qwen2.5-7B-Instruct), and one answer prompt,
over five question sets (NaturalQuestions, TriviaQA, HotpotQA,
2WikiMultihopQA, MuSiQue) totalling more than 37{,}000 questions.
Differences are tested with paired statistics (McNemar's test on Exact
Match, the Wilcoxon signed-rank test on continuous metrics) with
Holm-adjusted p-values, and every difference carries a paired bootstrap
95\% confidence interval. Budget-matched controls separate a better
architecture from more text in the context. The main results are:
(a)~cross-encoder reranking gains $+0.051$ Exact Match on HotpotQA and
loses $0.100$ on NaturalQuestions; (b)~the distribution of five retrieval
failure causes, measured per question, explains these reversals;
(c)~CascadeRAG wins \ResultWins{} of \ResultPairs{} paired comparisons
against every baseline on every question set, loses \ResultLosses{}, and
ties \ResultTies{}, at 1.16--1.58 generator calls per question;
(d)~the losslessness of the abstention trigger holds with zero exceptions
over all 887{,}917 (question, architecture) pairs measured; (e)~the result
replicates under a second, unrelated generator (Llama-3.1-8B-Instruct);
(f)~ablations isolate what the fused first stage and the escalation each
add, and show that a confidence threshold adds nothing. Every number in
the manuscript is generated by the benchmark code from the per-question
logs and imported as a named macro; the code, the logs, and the manuscript
pipeline are publicly released.

\section*{3. What papers by other authors make the most closely related
contributions, and how is your paper related to them?}

\begin{itemize}
\item \citet{jin2024flashrag} (FlashRAG) standardise datasets, corpora,
and reimplementations of a dozen RAG methods, and are the source of our
corpus and question sets. We add paired significance testing with
multiplicity control, the corpus-imposed ceiling, abstention accounting,
and a per-question failure diagnosis, and we study a different set of
architectures under one stack.
\item \citet{jeong2024adaptiverag} (Adaptive-RAG), \citet{asai2024selfrag}
(Self-RAG), and \citet{yan2024crag} (CRAG) all decide how much work a
question needs, from the question alone or with a trained component.
CascadeRAG decides \emph{after} the reader has seen the evidence (a
posterior signal) and requires no trained component; Adaptive-RAG and CRAG
are rows in our comparison.
\item \citet{chow1970reject} and \citet{geifman2017selective} formalise
the reject option and selective prediction; we show that the reader's
abstention is a reject option with zero selective risk under Exact Match
and exploit it as a deferral trigger.
\item \citet{viola2001cascade} and \citet{wang2011cascade} introduce
cascades in vision and in cascade ranking for retrieval; we apply the idea
at the level of whole RAG strategies.
\item \citet{wang2024bestpractices} tune RAG components in isolation; our
controlled comparison shows that the gains of such components depend on
the question distribution and can reverse.
\end{itemize}

\section*{4. Have you published parts of your paper before, for instance
in a conference?}

No. No part of this manuscript has been published before, in a conference
or elsewhere, and the manuscript is not under review at any other venue.

\bibliographystyle{apacite}
{\small\bibliography{references}}

\end{document}
